%% =====================================================================
%% Response-to-Reviewers Letter Template
%% --------------------------------------------------------------------
%% For use with IEEE TCAD / TVLSI / TODAES / ESL revisions.
%% This is a standalone document, compiled separately from the revised
%% manuscript. Submit both via ScholarOne (revised manuscript +
%% response letter + optional latexdiff PDF).
%% --------------------------------------------------------------------
%% Compile with: latexmk -pdf response.tex   or   pdflatex response.tex
%% =====================================================================

\documentclass[11pt,letterpaper]{article}

\usepackage[margin=1in]{geometry}
\usepackage{xcolor}
\usepackage{enumitem}
\usepackage{tcolorbox}
\usepackage{hyperref}
\hypersetup{hidelinks}
\usepackage[T1]{fontenc}
\usepackage{lmodern}
\usepackage{microtype}

%% ---- Colors ---------------------------------------------------------
\definecolor{commentbg}{RGB}{240,240,240}
\definecolor{commentframe}{RGB}{120,120,120}
\definecolor{responseaccent}{RGB}{0,90,140}
\definecolor{revexcerpt}{RGB}{0,0,160}

%% ---- Macros ---------------------------------------------------------

% \reviewercomment{ID}{verbatim quoted comment text}
% renders the quoted comment in a soft-gray box with the ID as a tag.
\newtcolorbox{reviewercomment}[1]{%
    colback=commentbg, colframe=commentframe,
    boxrule=0.4pt, arc=2pt, left=8pt, right=8pt, top=4pt, bottom=4pt,
    title={\textbf{#1}}, fonttitle=\small\sffamily,
    coltitle=black, attach boxed title to top left={xshift=4pt,yshift=-2pt},
    boxed title style={colback=commentbg, colframe=commentbg,
                       boxrule=0pt, left=4pt, right=4pt}}

% \response opens the response paragraph
\newcommand{\response}{\par\textbf{\color{responseaccent}Response:}\par}

% \paperref{Sec, p. X, lines Y-Z}
\newcommand{\paperref}[1]{\textit{(see #1 of the revised manuscript)}}

% \revexcerpt{...} — quote a piece of revised manuscript text
\newtcolorbox{revexcerpt}{%
    colback=white, colframe=revexcerpt,
    boxrule=0.4pt, arc=1pt, left=10pt, right=10pt,
    title={Excerpt from the revised manuscript}, fonttitle=\small\sffamily,
    coltitle=revexcerpt, attach boxed title to top left={xshift=4pt,yshift=-2pt},
    boxed title style={colback=white, colframe=white, boxrule=0pt}}

\title{\vspace{-1.5em}Response to Reviewers \\
       \large Manuscript ID TCAD-XXXX-XXXX \\
       \textit{[REPLACE WITH PAPER TITLE]}}
\author{}
\date{\today}

%% =====================================================================
\begin{document}
\maketitle
\vspace{-2em}

%% ---- Opening paragraph ----------------------------------------------

\noindent\textbf{To the Editor and Reviewers of IEEE TCAD:}

We sincerely thank the editor and the [\textbf{N}] reviewers for their
detailed and constructive comments on our manuscript. The reviews have
substantially strengthened the paper.

\paragraph{Main changes in this revision.}
\begin{itemize}[itemsep=2pt,topsep=2pt]
    \item \textbf{[REPLACE]} Added a new subsection (Sec.~III-D) providing
          a formal complexity analysis (R1.C2, R3.C1).
    \item \textbf{[REPLACE]} Added a new experimental subsection (Sec.~V-C)
          comparing against two additional recent baselines on the EPFL
          benchmark suite (R2.C1, R2.C4).
    \item \textbf{[REPLACE]} Substantially revised Sec.~II to better
          contextualize our contribution within prior work on
          [topic] (R1.C1, R2.C2).
    \item \textbf{[REPLACE]} Corrected the algorithmic-complexity
          discussion in Sec.~III-C (R1.C3).
    \item \textbf{[REPLACE]} Fixed all typos and notation inconsistencies
          flagged by Reviewer 3.
\end{itemize}

\noindent\textbf{Document layout.} Below we respond point-by-point to
the editor's letter and each reviewer's comments. Reviewer comments are
reproduced verbatim in shaded boxes; our responses follow in regular
type. References to the revised manuscript use the section / page / line
numbers of the revised version, with changes marked in blue text.

\hrule

%% =====================================================================
\section*{Response to the Editor}
%% =====================================================================

\begin{reviewercomment}{Editor's Letter}
\textit{[REPLACE: Quote the editor's specific guidance verbatim, e.g.,
"Reviewers raised concerns about the experimental coverage. In
particular, the lack of comparison to recent neural-network-based
approaches must be addressed in the revision."]}
\end{reviewercomment}

\response
We thank the editor for the clear summary of the deciding issues. We
have addressed each of them as described below; in particular:

\begin{itemize}[itemsep=2pt,topsep=2pt]
    \item \textbf{[Issue 1]:} Addressed in our response to R1.C2 and
          R2.C1, and reflected in the new Sec.~V-C
          \paperref{Sec.~V-C, p.~10, lines 12--38}.
    \item \textbf{[Issue 2]:} Addressed by the revisions described in
          our response to R1.C1.
    \item \textbf{[Issue 3]:} Addressed by the new Theorem~1 and the
          revised Sec.~III-D.
\end{itemize}

%% =====================================================================
\section*{Response to Reviewer 1}
%% =====================================================================

\paragraph{Summary.} We thank Reviewer~1 for the careful reading and the
constructive suggestions. \textbf{[REPLACE with a one-sentence summary
of the most impactful change you made for this reviewer.]}

\subsection*{R1.C1 --- [Short tag for the comment topic]}

\begin{reviewercomment}{Reviewer 1, Comment 1}
\textit{[REPLACE: Quote the reviewer's comment verbatim.]}
\end{reviewercomment}

\response We agree with the reviewer. \textbf{[REPLACE]} The previous
manuscript was unclear on this point. We have:

\begin{itemize}[itemsep=2pt,topsep=2pt]
    \item Added a new paragraph in Sec.~II-B
          \paperref{Sec.~II-B, p.~4, lines 8--19} explicitly delineating
          [the boundary].
    \item Cited [paper~X] and [paper~Y] as representative of the prior
          line of work the reviewer alluded to.
    \item Updated the contribution bullets in Sec.~I to clarify the
          scope.
\end{itemize}

\noindent For the reviewer's convenience, the relevant new text reads:

\begin{revexcerpt}
\textit{[REPLACE: Paste the new text from the revised manuscript
verbatim, in the same color used to mark revisions, e.g., blue.]}
\end{revexcerpt}

We hope this resolves the reviewer's concern.

\subsection*{R1.C2 --- [Short tag]}

\begin{reviewercomment}{Reviewer 1, Comment 2}
\textit{[REPLACE: Quote the reviewer's comment.]}
\end{reviewercomment}

\response We thank the reviewer for this important suggestion.
\textbf{[REPLACE]} In response, we have added a new subsection
(Sec.~III-D) providing a formal complexity analysis. Specifically:

\begin{itemize}[itemsep=2pt,topsep=2pt]
    \item Theorem~1 \paperref{Sec.~III-D, p.~7} now states the
          \(O(n \log n)\) complexity of the proposed algorithm.
    \item The proof is sketched in-line and given in full in
          Appendix~A \paperref{p.~16}.
    \item Figure~8 \paperref{Sec.~V-D, p.~12} now reports empirical
          scaling on designs from 10\,k to 500\,k cells, confirming
          the analytical result.
\end{itemize}

\subsection*{R1.C3 --- [Short tag, e.g., correction the reviewer flagged]}

\begin{reviewercomment}{Reviewer 1, Comment 3}
\textit{[REPLACE: Quote the reviewer's comment.]}
\end{reviewercomment}

\response The reviewer is correct, and we apologize for the original
oversight. \textbf{[REPLACE]} We have:

\begin{itemize}[itemsep=2pt,topsep=2pt]
    \item Corrected the complexity expression in Eq.~(5)
          \paperref{Sec.~III-C, p.~6, line 14}.
    \item Updated the discussion in the surrounding paragraph
          accordingly.
    \item Added a self-check experiment in Sec.~V-D verifying the
          corrected expression.
\end{itemize}

\subsection*{R1.C4, R1.C5, \ldots\ (typos and minor)}

We have addressed each of Reviewer~1's minor comments and typos in the
revised manuscript:

\begin{itemize}[itemsep=2pt,topsep=2pt]
    \item \textbf{[REPLACE: brief description of each]} (\paperref{location})
    \item \ldots
\end{itemize}

%% =====================================================================
\section*{Response to Reviewer 2}
%% =====================================================================

\paragraph{Summary.} We thank Reviewer~2 for the careful evaluation and
constructive comments.

\subsection*{R2.C1 --- [Short tag]}

\begin{reviewercomment}{Reviewer 2, Comment 1}
\textit{[REPLACE: Quote the reviewer's comment.]}
\end{reviewercomment}

\response We agree. \textbf{[REPLACE]} \ldots

\subsection*{R2.C2 --- (Respectful Disagreement Example)}

\begin{reviewercomment}{Reviewer 2, Comment 2}
\textit{[REPLACE: Quote a comment where you respectfully disagree.]}
\end{reviewercomment}

\response We thank the reviewer for raising this point. After
investigation, we believe the comparison the reviewer suggests is not
appropriate for the current submission, and we explain our reasoning
below.

\textbf{[REPLACE: Provide the justification, with citations to relevant
prior work. Be specific about scope, target technology, design level,
etc.]}

To address the reviewer's underlying concern, we have nonetheless added
a new paragraph in Sec.~II \paperref{Sec.~II-B, p.~4, lines 22--34}
that clarifies the boundary between [our scope] and [the scope suggested
by the reviewer], explicitly citing [paper~X] as representative of the
latter.

We hope this clarification resolves the concern. We are happy to add a
direct empirical comparison in a future extension if [the prerequisite]
becomes feasible.

%% =====================================================================
\section*{Response to Reviewer 3}
%% =====================================================================

\paragraph{Summary.} \textbf{[REPLACE with summary of changes for R3.]}

\subsection*{R3.C1 --- [Short tag]}

\begin{reviewercomment}{Reviewer 3, Comment 1}
\textit{[REPLACE: Quote.]}
\end{reviewercomment}

\response \textbf{[REPLACE.]}

\hrule

%% =====================================================================
\section*{Summary of Changes}
%% =====================================================================

For the reviewers' convenience, the major changes in this revision are
listed in one place below:

\begin{enumerate}[label=\textbf{Change \arabic*.}, leftmargin=*, itemsep=4pt]
    \item \textbf{[REPLACE]} New Sec.~III-D: formal complexity analysis
          (R1.C2, R3.C1). \emph{p.~7, lines 1--22.}
    \item \textbf{[REPLACE]} New Sec.~V-C: comparison to recent
          neural-network-based baselines on EPFL (R2.C1, R2.C4).
          \emph{p.~10, lines 12--38.}
    \item \textbf{[REPLACE]} Revised Sec.~II: clearer scope distinction
          and additional related work (R1.C1, R2.C2). \emph{pp.~3--5.}
    \item \textbf{[REPLACE]} Corrected complexity expression in
          Eq.~(5), with self-check experiment (R1.C3). \emph{p.~6,
          line 14, and Sec.~V-D.}
    \item \textbf{[REPLACE]} Typo fixes and notation cleanup throughout.
\end{enumerate}

\paragraph{New citations introduced in this revision.} All new citations
have been programmatically verified through IEEE Xplore, DBLP, and
CrossRef (see Sec.~Verification at the end of this document for
details). The new entries are: [paper X], [paper Y], [paper Z].

\paragraph{Diff PDF.} We have also submitted a \texttt{latexdiff}-generated
diff PDF (\texttt{main-diff.pdf}) as supplementary material to make the
changes easier to inspect.

\vspace{1em}

\noindent
We thank the editor and the reviewers again for the careful and
constructive review. We believe the revisions have substantially
strengthened the paper, and we look forward to your feedback.

\vspace{2em}

\noindent
Sincerely,\\
The Authors

\end{document}
